\chapter{Prawdopodobieństwo klasyczne, warunkowanie i niezależność}
\label{ch:conditioning}
\paragraph{Odsyłacze międzytomowe.} Zastosowania warunkowania do rozkładów a posteriori znajdują się w tomie IV, rozdziale~\ref{IV-ch:bayes-foundations}; modele sprzężone w tomie IV, rozdziale~\ref{IV-ch:conjugate-families}.


\section{Model Laplace'a}
Najprostszy pomysł: jeśli istnieje $N$ równie możliwych wyników,
a zdarzenie sprzyjające zawiera $k$ z nich, przypisujemy mu
szansę $k/N$. Jest to jednak \emph{warunkowe} zastosowanie
symetrii: trzeba najpierw uzasadnić, dlaczego wybrane wyniki
elementarne są równoprawdopodobne. Przy dwóch kostkach są nimi
uporządkowane pary oczek, nie sumy oczek. Wzór Laplace'a jest
konsekwencją wcześniejszej konstrukcji miary jednostajnej,
a nie dodatkowym aksjomatem.

\begin{definicja}
Modelem Laplace'a nazywamy przestrzeń probabilistyczną
$(\Omega,2^\Omega,\Prob)$, w której $\Omega$ jest niepustym zbiorem
skończonym, a $\Prob(\{\omega\})=1/|\Omega|$ dla każdego
$\omega\in\Omega$. Dla dowolnego $A\subseteq\Omega$ zachodzi wówczas
$\Prob(A)=|A|/|\Omega|$.
\end{definicja}

\begin{uwaga}
Warunek jednakowego prawdopodobieństwa wyników jest założeniem
modelu, a nie wnioskiem z samej symetrii językowego opisu.
Przykładowo podczas rzucania dwiema uczciwymi kostkami
36 uporządkowanych par ma równe prawdopodobieństwa; 11 możliwych sum
oczek nie ma rozkładu jednostajnego.
\end{uwaga}

\begin{przyklad}[Dwie kostki]
Niech $\Omega=\{1,\ldots,6\}^2$ z rozkładem jednostajnym.
Zdarzenie sumy równej siedem składa się z sześciu par,
więc ma prawdopodobieństwo $6/36=1/6$. Zdarzenie sumy równej dwa
składa się tylko z $(1,1)$, toteż jego prawdopodobieństwo wynosi $1/36$.
\end{przyklad}

\begin{przyklad}[Losowanie bez zwracania]
W urnie znajduje się $N$ różnych kul, z których $K$ jest czerwonych.
Losujemy jednocześnie $n$ kul, gdzie $0\leq n\leq N$, i zakładamy
jednostajność wszystkich $n$-elementowych podzbiorów.
Prawdopodobieństwo dokładnie $k$ czerwonych kul wynosi
\[
\frac{\binom Kk\binom{N-K}{n-k}}{\binom Nn}.
\]
W liczniku niezależnie \emph{zliczamy wybory} kul obu kolorów;
nie należy mylić tej kombinatorycznej zasady iloczynu
ze stochastyczną niezależnością kolejnych losowań.
\end{przyklad}

\section{Prawdopodobieństwo warunkowe}
Otrzymanie informacji zmienia przestrzeń rozważanych przypadków.
Jeśli wiemy, że zaszło $B$, to interesujące nas zdarzenie $A$
może zajść jedynie przez $A\cap B$. Aby pozostała pewna
możliwość miała prawdopodobieństwo jeden, normalizujemy
wagi przez $\Prob(B)$. To intuicja ilorazu
$\Prob(A\cap B)/\Prob(B)$.

Warunek $\Prob(B)>0$ nie jest formalnością: dla zdarzeń
zerowych dostalibyśmy iloraz $0/0$. Warunkowanie względem
wartości ciągłej zmiennej losowej wymaga później innej konstrukcji.

\begin{definicja}
Niech $(\Omega,\mathcal F,\Prob)$ będzie przestrzenią probabilistyczną,
a $A,B\in\mathcal F$ oraz $\Prob(B)>0$.
Prawdopodobieństwo warunkowe zdarzenia $A$ przy założeniu $B$
definiujemy wzorem
\begin{equation}\label{eq:cond}
\Prob(A\mid B):=\frac{\Prob(A\cap B)}{\Prob(B)}.
\end{equation}
\end{definicja}

\begin{stwierdzenie}
Dla każdego zdarzenia $B\in\mathcal F$ z $\Prob(B)>0$
funkcja $\Prob_B(A):=\Prob(A\mid B)$, $A\in\mathcal F$,
jest miarą probabilistyczną na $(\Omega,\mathcal F)$.
Ponadto $\Prob(B\mid B)=1$ i $\Prob(A^c\mid B)=1-\Prob(A\mid B)$.
\end{stwierdzenie}
\begin{proof}
Musimy sprawdzić wszystkie aksjomaty miary, nie tylko
nieujemność. Dla $A\in\mathcal F$ mamy
$0\leq\Prob(A\cap B)\leq\Prob(B)$, stąd
$\Prob_B(A)\in[0,1]$. Dalej
$\Prob_B(\Omega)=\Prob(\Omega\cap B)/\Prob(B)=1$.

Jeśli $A_1,A_2,\ldots$ są parami rozłączne, to zbiory
$A_n\cap B$ są również parami rozłączne, a
$B\cap\bigcup_n A_n=\bigcup_n(B\cap A_n)$.
Przeliczalna addytywność miary $\Prob$ daje
\[
\Prob_B\!\left(\bigcup_n A_n\right)
=\frac{\sum_n\Prob(A_n\cap B)}{\Prob(B)}
=\sum_n\Prob_B(A_n).
\]
Wreszcie $\Prob(B\mid B)=\Prob(B)/\Prob(B)=1$, zaś
$B=(A\cap B)\mathbin{\dot\cup}(A^c\cap B)$.
Dzieląc równość dla mas obu części przez $\Prob(B)$,
otrzymujemy $\Prob(A^c\mid B)=1-\Prob(A\mid B)$.
\end{proof}


\begin{przyklad}[Warunkowanie jako nowa miara na przestrzeni dwupunktowej]
Wróćmy do uczciwej monety
$\Omega=\{o,r\}$, $\Prob(\{o\})=\Prob(\{r\})=1/2$.
Warunkujemy względem zdarzenia $B=\{o\}$.
Wtedy $\Prob_B(\{o\})=1$ i $\Prob_B(\{r\})=0$.
Wynik doświadczenia nie zmienił się wstecznie;
zmieniła się miara, która opisuje naszą informację
po dowiedzeniu się, że zaszło $B$.
Ta prosta konstrukcja ilustruje, dlaczego rozkład warunkowy
jest nadal miarą probabilistyczną i dlaczego prawdopodobieństwa
$0$ oraz $1$ mogą pojawić się po warunkowaniu.
\end{przyklad}

\begin{przyklad}[Warunkowanie nie jest symetryczne]
W urnie znajdują się dwie kule czerwone i jedna niebieska.
Losujemy dwie kule kolejno bez zwracania. Niech $A$ oznacza,
że pierwsza jest czerwona, a $B$, że druga jest czerwona.
Z symetrii $\Prob(A)=\Prob(B)=2/3$.
Natomiast
$\Prob(A\cap B)=(2/3)(1/2)=1/3$, więc
$\Prob(B\mid A)=(1/3)/(2/3)=1/2$.
W tym szczególnym przykładzie również $\Prob(A\mid B)=1/2$,
co wynika z symetrii kolejności losowań; nie jest to jednak
ogólna reguła. Jeśli zdarzenia mają różne masy, na przykład
$\Prob(A)=1/2$, $\Prob(B)=1/4$ i
$\Prob(A\cap B)=1/8$, to
$\Prob(A\mid B)=1/2$, lecz $\Prob(B\mid A)=1/4$.
Dane te realizuje na przykład przestrzeń ośmiu
jednakowo prawdopodobnych wyników $\Omega=\{1,\ldots,8\}$,
w której $A=\{1,2,3,4\}$ oraz $B=\{1,5\}$.
\end{przyklad}

\begin{twierdzenie}[Reguła mnożenia]
Niech $A_1,\ldots,A_n\in\mathcal F$, gdzie $n\geq2$, i załóżmy,
że $\Prob(A_1\cap\cdots\cap A_{j})>0$ dla $1\leq j<n$.
Wtedy
\[
\Prob\!\left(\bigcap_{j=1}^n A_j\right)
=\Prob(A_1)\prod_{j=2}^n
\Prob\!\left(A_j\,\middle|\,\bigcap_{i=1}^{j-1}A_i\right).
\]
\end{twierdzenie}
\begin{proof}
Wprowadźmy $C_j=\bigcap_{i=1}^j A_i$ dla
$j=1,\ldots,n$. Założenia zapewniają,
że $\Prob(C_{j-1})>0$ dla $j=2,\ldots,n$;
wszystkie występujące prawdopodobieństwa warunkowe
są więc prawidłowo zdefiniowane. Ponieważ
$C_j=C_{j-1}\cap A_j$, z definicji mamy
\[
\Prob(A_j\mid C_{j-1})
=\frac{\Prob(C_j)}{\Prob(C_{j-1})}.
\]
Mnożąc te równości dla $j=2,\ldots,n$ i dopisując
$\Prob(A_1)=\Prob(C_1)$, otrzymujemy iloczyn
\[
\Prob(C_1)\frac{\Prob(C_2)}{\Prob(C_1)}
\frac{\Prob(C_3)}{\Prob(C_2)}\cdots
\frac{\Prob(C_n)}{\Prob(C_{n-1})}=\Prob(C_n).
\]
To właśnie prawdopodobieństwo przecięcia wszystkich $A_j$.
Na żadnym etapie nie założyliśmy niezależności:
każdy czynnik odpowiada sytuacji po uwzględnieniu
poprzednich zdarzeń.
\end{proof}

\begin{przyklad}[Bez zwracania i ze zwracaniem]
W urnie znajduje się pięć czerwonych i trzy niebieskie kule.
Przy dwóch losowaniach bez zwracania prawdopodobieństwo dwóch
czerwonych kul wynosi $(5/8)(4/7)=5/14$.
Przy losowaniu ze zwracaniem wynosi $(5/8)^2=25/64$.
Różnica bierze się stąd, że po pierwszym losowaniu bez zwracania
zmienia się skład urny, a więc rozkład wyniku drugiego losowania.
\end{przyklad}

\section{Niezależność zdarzeń}
W kontekście warunkowania można teraz powiedzieć precyzyjniej:
przy $\Prob(B)>0$ zdarzenie $B$ nie wnosi informacji o $A$
wtedy i tylko wtedy, gdy $\Prob(A\mid B)=\Prob(A)$.
Przykład trzech zdarzeń zależnych wyłącznie jako trójka
wyjaśni, dlaczego żądanie niezależności wszystkich skończonych
przekrojów jest silniejsze od niezależności parami.

Niezależność zdarzeń została zdefiniowana przed sformułowaniem
lematów Borela--Cantellego. Rozróżnienie niezależności wzajemnej
i parami niezależnej ma znaczenie również dla skończonych rodzin.


\begin{stwierdzenie}
Niech $A,B\in\mathcal F$ i $\Prob(B)>0$.
Wówczas $A$ i $B$ są niezależne wtedy i tylko wtedy, gdy
$\Prob(A\mid B)=\Prob(A)$. Jeżeli $A$ i $B$ są niezależne,
to niezależne są także pary $A^c,B$ oraz $A^c,B^c$.
\end{stwierdzenie}
\begin{proof}
Pierwsza równoważność wynika przez podstawienie do \eqref{eq:cond}.
Ponadto $\Prob(A^c\cap B)=\Prob(B)-\Prob(A\cap B)
=(1-\Prob(A))\Prob(B)$, a więc $A^c$ i $B$ są niezależne.
Powtarzamy rozumowanie po zamianie $B$ na $B^c$.
\end{proof}

\begin{przyklad}[Niezależność parami nie wystarcza]
Przeprowadzamy dwa niezależne rzuty uczciwą monetą.
W przestrzeni $\Omega=\{0,1\}^2$ z prawdopodobieństwem jednostajnym
rozważamy $A=\{(1,0),(1,1)\}$,
$B=\{(0,1),(1,1)\}$ i $C=\{(0,0),(1,1)\}$.
Każde zdarzenie ma prawdopodobieństwo $1/2$, a każde przecięcie
dwóch różnych z nich ma prawdopodobieństwo $1/4$.
Są zatem parami niezależne. Jednak $\Prob(A\cap B\cap C)=1/4$,
podczas gdy $\Prob(A)\Prob(B)\Prob(C)=1/8$.
\end{przyklad}

\begin{uwaga}
Zdarzenia rozłączne o dodatnich prawdopodobieństwach \emph{nie są}
niezależne: ich przecięcie ma prawdopodobieństwo zero,
a iloczyn ich prawdopodobieństw jest dodatni.
\end{uwaga}

\section{Prawdopodobieństwo całkowite i wzór Bayesa}
Jeżeli doświadczenie przebiega dwustopniowo, na przykład najpierw
wybieramy urnę, a potem kulę, możemy obliczyć szansę zdarzenia
końcowego osobno dla każdego sposobu rozpoczęcia doświadczenia.
Wzór na prawdopodobieństwo całkowite jest sumowaniem tych
rozłącznych dróg. Wzór Bayesa jest natomiast czysto algebraicznym
odwróceniem warunku: znając szansę obserwacji przy poszczególnych
scenariuszach, obliczamy szansę scenariusza pod warunkiem obserwacji.
Nie wprowadzamy tutaj jeszcze odrębnego aparatu wnioskowania
bayesowskiego.

\begin{twierdzenie}[Wzór na prawdopodobieństwo całkowite]
Niech $(B_i)_{i\in I}$ będzie skończoną lub przeliczalną rodziną
parami rozłącznych zdarzeń z $\Prob(B_i)>0$ dla każdego $i\in I$
oraz $\Prob(\bigcup_{i\in I}B_i)=1$.
Wtedy dla każdego $A\in\mathcal F$
\begin{equation}\label{eq:totalprob}
\Prob(A)=\sum_{i\in I}\Prob(A\mid B_i)\Prob(B_i).
\end{equation}
\end{twierdzenie}
\begin{proof}
Niech $B=\bigcup_{i\in I}B_i$.
Z założenia $\Prob(B)=1$, a więc $\Prob(B^c)=0$.
Rozkładamy $A$ na rozłączne części
$A=(A\cap B)\mathbin{\dot\cup}(A\cap B^c)$.
Ponieważ $A\cap B^c\subseteq B^c$, monotoniczność
daje $\Prob(A\cap B^c)=0$ i
$\Prob(A)=\Prob(A\cap B)$.

W dalszym kroku korzystamy z rozłączności $B_i$:
$A\cap B=\bigcup_{i\in I}(A\cap B_i)$,
a zbiory $A\cap B_i$ są parami rozłączne. Zatem
\[
\Prob(A)=\sum_{i\in I}\Prob(A\cap B_i)
=\sum_{i\in I}\Prob(A\mid B_i)\Prob(B_i).
\]
Ostatnia równość to definicja prawdopodobieństwa
warunkowego zastosowana osobno do każdego $B_i$.
Jeśli $I$ jest nieskończony, użyliśmy przeliczalnej,
a nie tylko skończonej addytywności.
\end{proof}

\begin{twierdzenie}[Wzór Bayesa]
Niech $(B_i)_{i\in I}$ spełnia warunki poprzedniego twierdzenia,
a $A\in\mathcal F$ ma dodatnie prawdopodobieństwo.
Wtedy dla każdego $j\in I$
\begin{equation}\label{eq:bayes}
\Prob(B_j\mid A)
=\frac{\Prob(A\mid B_j)\Prob(B_j)}
{\sum_{i\in I}\Prob(A\mid B_i)\Prob(B_i)}.
\end{equation}
\end{twierdzenie}
\begin{proof}
Z definicji warunkowania, ponieważ $\Prob(A)>0$,
\[
\Prob(B_j\mid A)=\frac{\Prob(B_j\cap A)}{\Prob(A)}.
\]
Licznik możemy zapisać jako
$\Prob(A\mid B_j)\Prob(B_j)$, gdyż $\Prob(B_j)>0$.
Natomiast mianownik jest sumą wkładów rozłącznych
scenariuszy $B_i$ z poprzedniego twierdzenia:
$\Prob(A)=\sum_{i\in I}\Prob(A\mid B_i)\Prob(B_i)$.
Po podstawieniu obu równości otrzymujemy tezę.
Nie ma tu nowego aksjomatu ani nowej interpretacji
prawdopodobieństwa: wzór wynika bezpośrednio z
warunkowania i rozbicia przestrzeni na rozłączne przypadki.
\end{proof}

\begin{przyklad}[Test diagnostyczny]
Niech $D$ oznacza występowanie pewnej choroby, a $T$ dodatni wynik
testu. Przyjmijmy $\Prob(D)=0{,}01$,
$\Prob(T\mid D)=0{,}99$ i $\Prob(T\mid D^c)=0{,}05$.
Wtedy
\[
\Prob(D\mid T)
=\frac{0{,}99\cdot0{,}01}
{0{,}99\cdot0{,}01+0{,}05\cdot0{,}99}
=\frac16.
\]
Dodatni wynik nie oznacza więc prawdopodobieństwa choroby równego
$99\%$. Wartość predykcyjna wyniku zależy zarówno od jakości testu,
jak i częstości choroby w rozpatrywanej populacji.
\end{przyklad}

\begin{przyklad}[Trzy drzwi --- problem Monty'ego Halla]
Załóżmy, że nagroda znajduje się jednostajnie za jednymi z trzech drzwi.
Gracz wybiera drzwi, nie znając położenia nagrody. Prowadzący zna jej
położenie, zawsze otwiera inne drzwi, za którymi nie ma nagrody,
a gdy ma wybór, rozstrzyga go losowo i symetrycznie. Gracz może
pozostać przy swoim wyborze albo zamienić go na jedne pozostałe
zamknięte drzwi. Początkowy wybór był prawidłowy z prawdopodobieństwem
$1/3$; wtedy zamiana przegrywa. Był nieprawidłowy z prawdopodobieństwem
$2/3$; wtedy prowadzący musi wskazać drugie puste drzwi, a zamiana
wygrywa. Zatem strategia zamiany daje prawdopodobieństwo wygranej $2/3$.
Warunki dotyczące zachowania prowadzącego są częścią modelu;
po ich zmianie wynik może być inny.
\end{przyklad}

\begin{przyklad}[Wybór urny po obserwacji koloru]
Najpierw losujemy urnę: pierwszą z prawdopodobieństwem $1/3$,
drugą z prawdopodobieństwem $2/3$. Pierwsza zawiera trzy kule
czerwone i jedną niebieską, druga jedną czerwoną i trzy niebieskie.
Następnie losujemy jednostajnie jedną kulę z wybranej urny.
Niech $U_1$ oznacza wybranie pierwszej urny, a $R$ wylosowanie
czerwonej kuli. Wtedy $\Prob(R\mid U_1)=3/4$ oraz
$\Prob(R\mid U_1^c)=1/4$. Po zobaczeniu czerwonej kuli
otrzymujemy ze wzoru Bayesa
\[
\Prob(U_1\mid R)
=\frac{(1/3)(3/4)}{(1/3)(3/4)+(2/3)(1/4)}
=\frac35.
\]
Przed obserwacją pierwsza urna była mniej prawdopodobna;
po obserwacji czerwonej kuli prawdopodobieństwo jej wyboru wzrasta.
Jest to wyłącznie rachunek prawdopodobieństwa warunkowego.
\end{przyklad}

\section{Warunkowanie a pozorna zależność}
Zależności widoczne w poszczególnych podgrupach mogą zniknąć
albo nawet zmienić znak po złączeniu obserwacji. Wynika to z tego,
że średnie warunkowe ważymy różnymi prawdopodobieństwami podgrup.
Tak zwany paradoks Simpsona nie jest paradoksem logicznym:
wskazuje, że przy porównywaniu danych trzeba jawnie określić,
na jakiej populacji warunkujemy.

\begin{przyklad}[Paradoks Simpsona]
W dwóch grupach porównujemy skuteczność metod $A$ i $B$.
W grupie I metoda $A$ daje $9$ sukcesów na $10$ prób, a metoda $B$
$80$ sukcesów na $100$ prób; w grupie II metoda $A$ daje
$20$ sukcesów na $100$ prób, a metoda $B$ $1$ sukces na $10$ prób.
W każdej grupie metoda $A$ ma większą częstość sukcesu:
$9/10>80/100$ oraz $20/100>1/10$.
Po połączeniu grup metoda $A$ ma jednak $29/110$ sukcesów,
a metoda $B$ $81/110$, więc kierunek porównania się odwraca.
Zmieniły się proporcje obu grup w populacjach stosujących dane metody.
Nie jest to sprzeczność, lecz ostrzeżenie przed utożsamianiem
prawdopodobieństw warunkowych i bezwarunkowych.
\end{przyklad}

